\documentclass[11pt]{article}
\usepackage[utf8]{inputenc}
\usepackage{amsmath, amssymb, amsfonts, amsthm}
\usepackage{geometry}
\geometry{a4paper, margin=1in}

\title{Stochastic Gradient Descent and the Langevin Equation}
\author{Analysis of Continuous-Time Asymptotics}
\date{\today}

\begin{document}

\maketitle

\section{Introduction}
Stochastic Gradient Descent (SGD) is a discrete-time optimization algorithm widely used in machine learning. Under the regime of a small learning rate $\eta$, its dynamics can be accurately modeled by a continuous-time stochastic differential equation (SDE), specifically the Overdamped Langevin Equation.

\section{The SGD Update Mechanism}
Consider an objective function $J(\theta)$ which is the mean of individual loss functions $J_i(\theta)$ over a dataset of size $n$:
\begin{equation}
J(\theta) = \frac{1}{n} \sum_{i=1}^{n} J_i(\theta)
\end{equation}

The standard SGD update at step $k$ for a randomly sampled index $i \in \{1, \dots, n\}$ is:
\begin{equation}
\theta_{k+1} = \theta_k - \eta \nabla J_i(\theta_k)
\end{equation}

\section{Decomposition of the Gradient}
To understand the relationship between SGD and the true gradient descent path, we decompose the stochastic gradient into the full gradient and a noise term:
\begin{equation}
\nabla J_i(\theta_k) = \nabla J(\theta_k) + \xi_k
\end{equation}
where $\xi_k = \nabla J_i(\theta_k) - \nabla J(\theta_k)$. By definition, the noise term has zero mean, $\mathbb{E}[\xi_k | \theta_k] = 0$.

The update rule can be rewritten as:
\begin{equation}
\theta_{k+1} - \theta_k = - \eta \nabla J(\theta_k) - \eta \xi_k
\end{equation}

\section{Continuous-Time Asymptotic (Langevin Limit)}
In the limit where $\eta \to 0$, we define a continuous-time process $\theta_t$ such that $\theta_k \approx \theta_{k\eta}$. Treating the noise term as a diffusion process, the trajectory follows the Langevin equation:
\begin{equation}
d\theta_t = -\nabla J(\theta_t) dt + \sqrt{2\beta^{-1}} dW_t
\end{equation}
where $dW_t$ represents a Wiener process and $\beta$ is the inverse temperature parameter related to the learning rate and the covariance of the gradient noise $\xi$.

\section{Significance}
This mapping allows for several theoretical insights:
\begin{itemize}
    \item \textbf{Exploration:} The diffusion term allows $\theta$ to explore the landscape and escape local minima.
    \item \textbf{Stationary Distribution:} The process converges to the Gibbs distribution $p(\theta) \propto e^{-\beta J(\theta)}$.
    \item \textbf{Stability:} The continuous form makes it easier to apply tools from calculus and PDE theory to study convergence rates.
\end{itemize}

\end{document}